% Part: sets-functions-relations
% Chapter: relations
% Section: special-properties

\documentclass[../../../include/open-logic-section]{subfiles}

\begin{document}

\olfileid{sfr}{rel}{prp}
\olsection{关系的特殊性质}

\begin{intro}
有些种类的关系非常常见，以至于人们为它们赋予了专门的名称。例如，$\le$ 和~$\subseteq$ 都以相似的方式关联着它们各自的定义域（对于~$\le$ 是~$\Nat$，对于~$\subseteq$ 是~$\Pow{A}$）。为了精确地说明这些关系何以相似，又何以不同，我们根据关系可能具有的一些特殊性质来对它们进行分类。事实上，这些特殊性质中的某些（组合）尤为重要：序关系和等价关系。
\end{intro}

\begin{defn}[自反性]
关系 $R \subseteq A^2$ 是\emph{自反的}，当且仅当对于每个 $x \in A$，都有 $Rxx$。
\end{defn}

\begin{defn}[传递性]
关系 $R \subseteq A^2$ 是\emph{传递的}，当且仅当只要 $Rxy$ 且 $Ryz$，则也有 $Rxz$。
\end{defn}

\begin{defn}[对称性]
关系~$R \subseteq A^2$ 是\emph{对称的}，当且仅当只要 $Rxy$，则也有~$Ryx$。
\end{defn}

\begin{defn}[反对称性]
关系~$R \subseteq A^2$ 是\emph{反对称的}，当且仅当只要 $Rxy$ 和 $Ryx$ 同时成立，就有 $x=y$（换言之：如果 $x\neq y$，则要么 $\lnot Rxy$，要么 $\lnot Ryx$）。
\end{defn}

\begin{explain}
在对称关系中，$Rxy$ 和 $Ryx$ 总是同时成立，或者同时不成立。在反对称关系中，$Rxy$ 和 $Ryx$ 能同时成立的唯一方式是 $x = y$。注意，这并不\emph{要求}当 $x = y$ 时 $Rxy$ 和 $Ryx$ 成立，而只是说这不被排除。因此，一个反对称关系可以是自反的，但并非每个反对称关系都是自反的。还要注意，“反对称”与仅仅是“不对称”是不同的条件。事实上，一个关系可以同时是对称且反对称的（例如，恒等关系就是如此）。
\end{explain}

\begin{defn}[连通性]
关系 $R \subseteq A^2$ 是\emph{连通的}，如果对于所有 $x,y\in A$，只要 $x \neq y$，则要么 $Rxy$，要么~$Ryx$。
\end{defn}

\begin{prob}
请举出满足以下条件的关系的例子： 自反且对称，但不传递； 自反且反对称； 反对称、传递，但不自反； 自反、对称且传递。不要使用关于数或集合的关系。
\end{prob}

\begin{defn}[反自反性]
关系 $R \subseteq A^2$ 被称为\emph{反自反的}，如果对于所有 $x \in A$，$Rxx$ 均不成立。
\end{defn}

\begin{defn}[非对称性]
关系 $R \subseteq A^2$ 被称为\emph{非对称的}，如果不存在一对 $x,y\in A$ 使得 $Rxy$ 和~$Ryx$ 同时成立。
\end{defn}

注意，如果 $A \neq \emptyset$，那么 $A$ 上任何反自反关系都不是自反的，并且 $A$ 上任何非对称关系也都是反对称的。然而，存在既不自反也不反自反的 $R \subseteq A^2$，也存在是反对称但不是非对称的关系。

\end{document}